\section{Conclusion and Future Work}
\label{sec:conclusion}

\subsection{Findings}
Three findings stand out from the analysis in
Section~\ref{sec:sql}:
\begin{enumerate}
  \item \textbf{Volume is not coverage.} Exfiltration labels dominate
    the corpus by raw line count (roughly 53{,}000 labeled lines),
    but this is an artefact of the DNS-tunneling phase emitting one
    labeled line per DNS query. The privilege-escalation phase has
    only 82 labeled lines but spans five distinct log sources, and
    the operational severity of an account takeover is incomparably
    higher. A detection program optimised for label volume would
    miss the most consequential phase entirely.
  \item \textbf{Cross-source joins are diagnostic.} Queries A4 and
    A6, plus the saved view
    \texttt{v\_privilege\_escalation\_timeline}, reconstruct the
    privilege-escalation sequence and the cross-host lateral movement
    using foreign-key joins between audit events and attack labels.
    Neither view is recoverable from any single log file in
    isolation; the database makes both legible in a single ordering.
  \item \textbf{Timing-only detection is feasible for short bursts.}
    Query A7 uses only \texttt{LAG} over event timestamps and detects
    two distinct attack bursts on \texttt{intranet\_server} without
    consulting a single label. The two-burst structure recovered from
    timing alone matches the structure recovered from labels in
    query A4 exactly. This is a useful unsupervised signal that
    complements rule-based detection.
\end{enumerate}

\subsection{What Happened Next: Exfiltration on \texttt{internal\_share}}
The deck and the report focus on the
\texttt{initial\_access}-into-\texttt{privilege\_escalation} narrative
on \texttt{intranet\_server}, but the dataset's full attack chain
continues for several more hours and migrates to a second host. At
13:50 UTC on the same day (about nine hours after the
\texttt{sudo cat /etc/shadow} burst), \texttt{internal\_share} emits
two audit events for a service called \texttt{put} that starts and
stops within seconds. Query A6 surfaces both events in the same
ordering as the privilege-escalation chain, and the basic
service-activity baseline (B6) confirms that the \texttt{put} service
has no other appearances in the audit data. In the labels domain, the
DNS-tunneling rule \texttt{dnsteal.domain.match} fires roughly
53{,}000 times on the same window, indicating that
\texttt{/etc/shadow} or downstream credentials material was
exfiltrated through DNS queries to the attacker-controlled domain.
The report does not deep-dive the exfiltration phase further: the
loaded slice only contains the audit-domain side of the
exfiltration story, and the more interesting DNS-tunneling channel
lives in a log type that future work will load into the schema (see
Future Work below). The exfiltration phase appears in the dashboard's
KPI band when the user widens the host filter to include
\texttt{internal\_share} and is reproduced in the appendix.

\subsection{Limitations}
\begin{itemize}
  \item \textbf{Single testbed slice.} The current iteration covers
    the \texttt{russellmitchell} testbed only. The other four
    testbeds in AIT-LDSv2.0 share the same schema and could be
    appended with no model changes; the work is loader-only.
  \item \textbf{Out-of-scope log domains.} DNS, OpenVPN, Apache, and
    system auth logs were explored in notebooks but are not loaded
    into the relational schema. Bringing them in would lengthen the
    SQL portfolio (especially for the exfiltration phase) but does
    not change the schema design. The DNS-tunneling channel, in
    particular, is the natural next domain to add: it is the
    highest-line-count phase of the attack and the report currently
    only describes it via the labels domain.
  \item \textbf{Single-machine deployment.} The PostgreSQL container
    runs on one developer machine. The current corpus is small enough
    that this is fine; scaling to the full AIT-LDSv2.0 with all
    testbeds and all log domains would push past a single instance
    and invite the distributed-SQL discussion.
  \item \textbf{Static dashboard scope.} The Streamlit dashboard
    currently shows aggregate views, a filterable timeline, and a
    static EXPLAIN summary. A real SOC dashboard would also support
    time-window scrubbing, alert drill-down, and per-rule
    false-positive review. We treat these as future work rather than
    as gaps.
  \item \textbf{Story-canonical defaults.} The story panels are
    pinned to the \texttt{intranet\_server} narrative so the
    audience does not get lost while experimenting with filters.
    This keeps the dashboard reproducible for grading but means a
    user who wants to investigate a different host has to read the
    SQL and run it manually rather than re-using the UI.
\end{itemize}

\subsection{Future Work}
Concrete and bounded next steps, in order of priority:
\begin{enumerate}
  \item \textbf{Load the remaining four testbeds.} The schema is
    testbed-agnostic; the loader needs a single configuration change
    and an additional run. Brings the labeled corpus from roughly
    62{,}000 to roughly 300{,}000 rows.
  \item \textbf{Add DNS and OpenVPN domains.} The DNS-tunneling
    exfiltration is severely under-represented in the audit-only
    schema. Adding \texttt{dns\_query} and \texttt{vpn\_session}
    tables (with obvious foreign keys to \texttt{host}) would let
    queries like A14 (detection-rule effectiveness) compare
    \texttt{dnsteal.*} rules against audit-based rules on equal
    footing, on the same population.
  \item \textbf{Materialize \texttt{v\_privilege\_escalation\_timeline}.}
    As the corpus grows, the underlying view will become a query-time
    bottleneck. Promoting it to a materialized view with an
    incremental-refresh schedule is a straightforward optimisation
    and provides natural snapshot semantics for the report.
  \item \textbf{Per-host story panels.} The current dashboard pins
    panels 1 to 8 to \texttt{intranet\_server}. Generalising the
    panel SQL so the host filter actually drives the story panels
    (with the privilege-escalation arc replaced by the relevant arc
    for whichever host the user picks) would make the dashboard a
    real investigation tool rather than a curated demo.
  \item \textbf{Productionize the dashboard.} Add user
    authentication, a time-range picker, and a Slack-friendly export
    of the attack timeline. None of these require schema changes.
\end{enumerate}
